\section{Notation and definitions}
\label{sec:notation}

This section defines every symbol used later.  A reader who knows quantum circuits can skim it;
a reader who does not should find here everything needed to follow the rest.

\paragraph{Qubits, states, measurement.}
A \emph{qubit} is a two-level quantum system with basis states $\ket{0}$ and $\ket{1}$.  The
state of $n$ qubits is a unit vector $\psi \in \mathbb{C}^{2^n}$ whose components are indexed by
$n$-bit strings; we write the string with \emph{qubit~1 (the first input) as the most significant
bit}, matching the paper's circuit figures, which number qubits from the top.  Measuring all
qubits returns the bit string $x$ with probability $|\psi_x|^2$.  When only some qubits are
measured, the probability of their outcome is the sum of $|\psi_x|^2$ over the strings $x$ that
agree with it on those qubits (the \emph{marginal}).  A \emph{mixed} state, needed as soon as
noise enters, is a positive matrix $\rho$ of unit trace (the \emph{density matrix}); a pure state
is $\rho=\psi\psi^\dagger$, and measurement probabilities are its diagonal entries.

\paragraph{Gates.}
A gate is a unitary matrix acting on one or two qubits.  The ones the paper uses are
\begin{align}
\RX(\theta) &= \begin{pmatrix}\cos\frac{\theta}{2} & -i\sin\frac{\theta}{2}\\[2pt]
-i\sin\frac{\theta}{2} & \cos\frac{\theta}{2}\end{pmatrix}, \qquad
H = \tfrac{1}{\sqrt2}\begin{pmatrix}1&1\\1&-1\end{pmatrix}, \\
\CRX(\theta) &= \ket{0}\bra{0}\otimes I + \ket{1}\bra{1}\otimes \RX(\theta)
\quad\text{(rotate the target by $\theta$ only if the control is $\ket{1}$).}
\end{align}
$\RX(\theta)\ket{0} = \cos\frac{\theta}{2}\ket{0} - i\sin\frac{\theta}{2}\ket{1}$, so after
$\RX(\theta)$ a qubit prepared in $\ket{0}$ is measured as $1$ with probability
$\sin^2(\theta/2)$.  $H$ turns $\ket{0}$ into the equal superposition
$(\ket{0}+\ket{1})/\sqrt2$; $H$ on every input qubit prepares all $2^n$ input strings at once with
equal weight.

\paragraph{Encoding a real number as an angle.}
A feature value $s\in[0,1]$ enters the circuit as $\RX(\theta(s))$ on its input qubit.  The paper
does not state $\theta(s)$; its companion paper \cite{snn} (Fig.~4) prints
$\CRX(1.369)$ for the value $0.4$ and $\CRX(1.7722)$ for $0.6$, which pins
\begin{equation}
\theta(s) = 2\arcsin\sqrt{s}, \qquad\text{so that } P(\text{qubit reads } 1) = s .
\label{eq:encoding}
\end{equation}
We use \cref{eq:encoding} throughout and report the alternative $\theta = \pi s$ as an ablation
where it matters (ECG, \cref{sec:ecg}).

\paragraph{Networks and their output.}
A network has $n_{\mathrm{in}}$ input qubits, $n_{\mathrm{hid}}$ hidden qubits (one per hidden
neuron) and one output qubit; $n = n_{\mathrm{in}}+n_{\mathrm{hid}}+1$.  Its output for a given
input is the \emph{activation}
\begin{equation}
\Pout \;=\; \Pr[\text{output qubit reads }1],
\end{equation}
a number in $[0,1]$.  For a classifier it is thresholded (Flappy: flap iff $\Pout>\tfrac12$);
for a regressor it is the predicted value (ECG: the denoised sample).

\paragraph{The superposition observable and the training loss.}
For the bit-count and Iris tasks the inputs are prepared in superposition and the measured joint
distribution of (input string $x$, output bit) is read.  The \emph{model distribution} is
\begin{equation}
q(x) \;=\; \frac{P(\text{inputs}=x,\ \mathrm{out}=1)}{\sum_{x'} P(\text{inputs}=x',\ \mathrm{out}=1)},
\label{eq:model-distribution}
\end{equation}
i.e.\ the input distribution \emph{conditioned on the output having fired} --- the output
amplitude acts as a weight on each input pattern (paper, Fig.~4 caption).  It is compared to a
target distribution $p(x)$ by the Kullback--Leibler divergence (paper, Eq.~(4)),
\begin{equation}
\KL(p\,\|\,q) \;=\; \sum_x p(x)\,\log\frac{p(x)}{q(x)} \;\ge 0,
\qquad \KL = 0 \iff p = q,
\label{eq:kl}
\end{equation}
with $0\log 0 = 0$ and $q$ floored at $10^{-12}$ so the value stays finite (our choice; the paper
does not say).  A second dissimilarity we use for noise is the \emph{total-variation distance}
$\TV(p,q) = \tfrac12\sum_x |p(x)-q(x)| \in [0,1]$, the largest difference in probability that any
event can have under the two distributions.

\paragraph{Noise channels.}
Real gates are not unitary.  A \emph{channel} is a linear map $\rho\mapsto\sum_i K_i\rho K_i^\dagger$
with $\sum_i K_i^\dagger K_i = I$ (the $K_i$ are its \emph{Kraus operators}); it sends density
matrices to density matrices.  Three channels describe a calibrated machine:
\begin{itemize}[nosep]
\item \emph{Depolarising} with probability $p$: with probability $1-p$ nothing happens, with
probability $p$ the state of the affected qubit(s) is replaced by the maximally mixed state.  A
vendor's ``average gate error'' $\varepsilon$ is converted to $p$ by
$\varepsilon = p\,(1 - 2^{-k})\,\tfrac{2^k}{2^k+1}$ for a $k$-qubit gate ($p = 2\varepsilon$ for
one qubit, $p = \tfrac{4}{3}\cdot\tfrac{5}{4}\varepsilon$ for two), which makes the channel's
average fidelity equal to $1-\varepsilon$ (cross-checked against Qiskit~Aer).
\item \emph{Relaxation} over a time $t$ with the vendor's $T_1$ (energy decay time) and $T_2$
(phase decay time): $\ket{1}$ decays toward $\ket{0}$ with probability $1-e^{-t/T_1}$ and
superpositions lose coherence by $e^{-t/T_2}$.  It is applied during every gate (for its
duration) and during every idle interval, which we schedule as early as possible, exactly as
IBM's compiler does.
\item \emph{Readout confusion}: the recorded bit differs from the true one with probability
$p_{01}$ (true $0$ read as $1$) or $p_{10}$ (true $1$ read as $0$).
\end{itemize}
The resulting \emph{noise model} of a machine is the rule ``after each gate, apply the gate's
depolarising channel and its relaxation; on every idle qubit, apply relaxation for the idle time;
at the end, apply readout confusion''.  Its parameters come from the vendor data of
\cref{app:provenance}; a \emph{noise scale} $s$ multiplies all of $p$, $1/T_1$, $1/T_2$ and the
readout errors together ($s=1$ is the published machine, $s=0$ is noiseless).

\paragraph{Shots.}
A quantum computer does not return probabilities; it returns one bit string per run, and
probabilities are estimated by repeating the circuit $N$ times (\emph{$N$ shots}).  An estimate
$\hat p = k/N$ of a probability $p$ has standard deviation $\sqrt{p(1-p)/N}\le 1/(2\sqrt N)$.  We
call $\delta = |\Pout - \tfrac12|$ the \emph{decision margin} of a thresholded output: a frame
decided from $N$ shots is decided wrongly with probability $\approx\Phi(-2\delta\sqrt N)$, where
$\Phi$ is the standard normal distribution function.

\paragraph{Fidelity.}
For a compiled circuit on a given machine we quote the estimate
\begin{equation}
F \;\approx\; \prod_{\text{2q gates}}(1-\varepsilon_{2})\;\prod_{\text{1q gates}}(1-\varepsilon_{1})\;
\prod_{\text{qubits}} F_{\mathrm{idle}}(t)\;\prod_{\text{qubits}}\bigl(1-\tfrac{p_{01}+p_{10}}{2}\bigr),
\label{eq:fidelity}
\end{equation}
the probability that no error at all occurred in the run, with
$F_{\mathrm{idle}}(t) = \tfrac12 + \tfrac16\bigl(2e^{-t/T_2}+e^{-t/T_1}\bigr)$ the average
fidelity of relaxation over the circuit's duration $t$ \cite{nielsen}.  It is a first-order
bookkeeping estimate, not a simulation; the simulations are reported next to it.

\paragraph{Abbreviations.}
KW--SA: the paper's optimiser (\cref{sec:kwsa}).  PSNR: peak signal-to-noise ratio,
$10\log_{10}(R^2/\mathrm{MSE})$ with $R$ the signal range and MSE the mean squared error, in
decibels (dB).  SSIM: structural similarity index \cite{wang}, a number in $[0,1]$ comparing
local means, variances and covariance of two signals.  SWAP: the gate that exchanges two qubits,
inserted by a compiler when a two-qubit gate is requested between qubits that are not physically
connected; it costs three native two-qubit gates.
